% GENERATED FILE. Edit canonical lessons, then run npm run generate:books.
% canonical-source-hash: fnv1a64:9aa9e7abfa07ce53
% canonical-lessons: SA-C128-bhayankarah, SA-C128-mithya, SA-C128-katuh, SA-C128-amlah, SA-C128-tiktah

\chapter{Terrible, False, Pungent, Sour, Bitter}
\label{ch:sa-terrible-false-pungent-sour-bitter}
\hlchaptermodality{128}

\begin{chapteropening}
\textbf{By the end of this chapter:} \emph{I can say terrible, false, pungent, sour and bitter.}

Complete the last lesson of chapter 128: i can say terrible, false, pungent, sour and bitter.
\end{chapteropening}

\section[\texorpdfstring{\sk{भयंकरः} (bhayaṅkaraḥ)}{भयंकरः (bhayaṅkaraḥ)}]{\sk{भयंकरः} (bhayaṅkaraḥ) — terrible}

\label{lesson:SA-C128-bhayankarah}



\begin{quote}
Before the new one: say the Sanskrit for famous, then the Sanskrit for safe.
\end{quote}

\subsection*{You'll want to know: \sk{भयंकरः}}
\textbf{\sk{भयंकरः}} — \emph{bhayaṅkaraḥ} — \textquotedblleft{}terrible\textquotedblright{}.

\begin{grammarlens}[title={What you've built}]
One more word for this chapter.
\end{grammarlens}

\subsection*{Your turn}
\begin{itemize}
\raggedright
  \item \emph{Say it:} \emph{bhayaṅkaraḥ}
  \item \emph{Say it:} \emph{bhayaṅkaraḥ}, once more
  \item \emph{Say it:} say \emph{bhayaṅkaraḥ}
  \item \emph{Recall:} say the Sanskrit for famous, then the Sanskrit for safe, then say \emph{bhayaṅkaraḥ} again
\end{itemize}

\begin{tcolorbox}[breakable,colback=teal!4,colframe=teal!35!black,title={Before you move on}]
What does \sk{भयंकरः} mean? (\textquotedblleft{}Terrible\textquotedblright{}.) Say it once more.
\end{tcolorbox}
\section[\texorpdfstring{\sk{मिथ्या} (mithyā)}{मिथ्या (mithyā)}]{\sk{मिथ्या} (mithyā) — false}

\label{lesson:SA-C128-mithya}



\begin{quote}
Before the new one: say the Sanskrit for safe, then the Sanskrit for terrible.
\end{quote}

\subsection*{You'll want to know: \sk{मिथ्या}}
\textbf{\sk{मिथ्या}} — \emph{mithyā} — \textquotedblleft{}false\textquotedblright{}.

\begin{grammarlens}[title={What you've built}]
One more word for this chapter.
\end{grammarlens}

\subsection*{Your turn}
\begin{itemize}
\raggedright
  \item \emph{Say it:} \emph{mithyā}
  \item \emph{Say it:} \emph{mithyā}, once more
  \item \emph{Say it:} say \emph{mithyā}
  \item \emph{Recall:} say the Sanskrit for safe, then the Sanskrit for terrible, then say \emph{mithyā} again
\end{itemize}

\begin{tcolorbox}[breakable,colback=teal!4,colframe=teal!35!black,title={Before you move on}]
What does \sk{मिथ्या} mean? (\textquotedblleft{}False\textquotedblright{}.) Say it once more.
\end{tcolorbox}
\section[\texorpdfstring{\sk{कटुः} (kaṭuḥ)}{कटुः (kaṭuḥ)}]{\sk{कटुः} (kaṭuḥ) — pungent}

\label{lesson:SA-C128-katuh}



\begin{quote}
Before the new one: say the Sanskrit for terrible, then the Sanskrit for false.
\end{quote}

\subsection*{You'll want to know: \sk{कटुः}}
\textbf{\sk{कटुः}} — \emph{kaṭuḥ} — \textquotedblleft{}pungent\textquotedblright{}.

\begin{grammarlens}[title={What you've built}]
One more word for this chapter.
\end{grammarlens}

\subsection*{Your turn}
\begin{itemize}
\raggedright
  \item \emph{Say it:} \emph{kaṭuḥ}
  \item \emph{Say it:} \emph{kaṭuḥ}, once more
  \item \emph{Say it:} say \emph{kaṭuḥ}
  \item \emph{Recall:} say the Sanskrit for terrible, then the Sanskrit for false, then say \emph{kaṭuḥ} again
\end{itemize}

\begin{tcolorbox}[breakable,colback=teal!4,colframe=teal!35!black,title={Before you move on}]
What does \sk{कटुः} mean? (\textquotedblleft{}Pungent\textquotedblright{}.) Say it once more.
\end{tcolorbox}
\section[\texorpdfstring{\sk{अम्लः} (amlaḥ)}{अम्लः (amlaḥ)}]{\sk{अम्लः} (amlaḥ) — sour}

\label{lesson:SA-C128-amlah}



\begin{quote}
Before the new one: say the Sanskrit for false, then the Sanskrit for pungent.
\end{quote}

\subsection*{You'll want to know: \sk{अम्लः}}
\textbf{\sk{अम्लः}} — \emph{amlaḥ} — \textquotedblleft{}sour\textquotedblright{}.

\begin{grammarlens}[title={What you've built}]
One more word for this chapter.
\end{grammarlens}

\subsection*{Your turn}
\begin{itemize}
\raggedright
  \item \emph{Say it:} \emph{amlaḥ}
  \item \emph{Say it:} \emph{amlaḥ}, once more
  \item \emph{Say it:} say \emph{amlaḥ}
  \item \emph{Recall:} say the Sanskrit for false, then the Sanskrit for pungent, then say \emph{amlaḥ} again
\end{itemize}

\begin{tcolorbox}[breakable,colback=teal!4,colframe=teal!35!black,title={Before you move on}]
What does \sk{अम्लः} mean? (\textquotedblleft{}Sour\textquotedblright{}.) Say it once more.
\end{tcolorbox}
\section[\texorpdfstring{\sk{तिक्तः} (tiktaḥ)}{तिक्तः (tiktaḥ)}]{\sk{तिक्तः} (tiktaḥ) — bitter}

\label{lesson:SA-C128-tiktah}



\begin{quote}
Before the new one: say the Sanskrit for pungent, then the Sanskrit for sour.
\end{quote}

\subsection*{You'll want to know: \sk{तिक्तः}}
\textbf{\sk{तिक्तः}} — \emph{tiktaḥ} — \textquotedblleft{}bitter\textquotedblright{}.

\begin{grammarlens}[title={What you've built}]
One more word for this chapter.
\end{grammarlens}

\subsection*{Your turn}
\begin{itemize}
\raggedright
  \item \emph{Say it:} \emph{tiktaḥ}
  \item \emph{Say it:} \emph{tiktaḥ}, once more
  \item \emph{Say it:} say \emph{tiktaḥ}
  \item \emph{Recall:} say the Sanskrit for pungent, then the Sanskrit for sour, then say \emph{tiktaḥ} again
\end{itemize}

\begin{tcolorbox}[breakable,colback=teal!4,colframe=teal!35!black,title={Before you move on}]
What does \sk{तिक्तः} mean? (\textquotedblleft{}Bitter\textquotedblright{}.) Say it once more.
\end{tcolorbox}
